\subsection{Deep Representation Learning}\label{subsec:mm-deep-learning}

Before calculating evolutionary disparities, it is necessary to project the raw visual semantics of organisms into a mathematically tractable morphological space. 
We define a mapping function $f_\theta$ parameterised by a Convolutional Neural Network (CNN). 
For a given image of an organism, the network directly processes the pixel matrix and extracts a high-dimensional feature vector from its final fully connected (fc) layer, bypassing the requirement for homologous anatomical coordinates.
% The extracted vector is subsequently $L_2$-normalised. Every organism is mathematically represented as a unit vector $v$. 
% This normalisation forces the deep visual morphospace to function geometrically as a unit hypersphere $\mathbb{S}^{D-1}$, 
% where $D$ is the dimensionality of the feature space. Within this Riemannian manifold, semantic and morphological similarity is encoded in angular directions rather than spatial magnitudes.

\subsection{Geometric Foundations of the Deep Morphospace}\label{subsec:mm-geometric-foundations}

To employ classical macroevolutionary models (such as Brownian motion and Phylogenetic Independent Contrasts) 
on feature representations extracted from deep neural networks, it is crucial to define the precise geometry of the underlying feature space. 
As demonstrated by \textcite{Wang2017NormFace} in the context of hyperspherical embedding, the loss function with $L_2$ normalization forces 
the network to minimize intra-class angular variance while maximizing inter-class angular disparity. 
Consequently, the semantic and morphological information learned by the model is primarily encoded in the angular direction of the feature vectors, 
rather than their Euclidean magnitudes.

Mathematically, the $L_2$ normalisation maps any vector $\mathbf{x} \in \mathbb{R}^D$ onto a unit hypersphere $\mathbb{S}^{D-1}$, defined as:

$$\mathbf{z} = \frac{\mathbf{x}}{\|\mathbf{x}\|_2}$$

where $\|\mathbf{z}\|_2 = 1$. 
Therefore, the deep visual morphospace in our framework does not function as a flat Euclidean space, 
but rather as a curved Riemannian manifold -- specifically, a unit hypersphere. 
The phenotypic distance between any two taxa $i$ and $j$ cannot be measured by Euclidean distance, which would violate the constraints of the surface topology. 
Instead, it is quantified via the geodesic distance (the great-circle distance) or the angular separation $\theta_{ij}$, formulated as:

$$\theta_{ij} = \arccos(\mathbf{z}_i \cdot \mathbf{z}_j)$$

This hyperspherical formulation provides the theoretical and geometric justification for our subsequent developments in spherical 
Ancestral State Reconstruction (ASR) and Riemannian Phylogenetic Independent Contrasts (PIC). 

\subsection{Similarity and Disparity}\label{subsec:mm-similarity-clustering}

According to the aforementioned geometric foundations, species similarity can be quantified using cosine similarity, implemented via
dot product calculation of $L_2$-normalised vectors in PyTorch \parencite{Ansel2024Pytorch}.
Next, we conducted agglomerative hierarchical clustering using the average
linkage method to merge clusters. This was executed with the
hierarchical clustering functionalities implemented in the SciPy
library \parencite{Gommers2025Scipy}. The hierarchical structure of the
clusters was output in Newick format, a widely used format in
computational biology for tree structures. Finally, we utilised ETE3
to export the clustering dendrogram in SVG format \parencite{Huerta2016ETE}.

To analyse the clustering result, we applied a recursive top-down
analysis to the tree, evaluating each internal node for taxonomic
"purity." For a given node, we defined taxonomic purity as
the proportion of the majority taxon. A node with a taxonomic
purity of more than 85\% was considered taxonomically consistent. For
taxonomically consistent nodes, all of their child nodes were
excluded from further checks. For such nodes, we further examined
all species to identify and annotate taxonomical outliers, which belong
to taxa that different from the majority taxon of this branch.
Finally, we conducted manual review to confirm whether outliers had
biological similarity with the majority taxa of their branches and
what kinds of similarity do they have. 
% The above pipeline was carried out in both order-level and family-level.

Similarly, spherical vector variance of species is utilised to assess the morphological
disparity among taxa. In the case study, we will examine the reliability of using spherical variance as a measure of morphological disparity.

\subsection{Disparity through time}\label{subsec:mm-disparity-through-time}

In classical phylogenetic independent contrast (PIC) proposed by
\textcite{Felsenstein1985PIC}, The evolution of phenotypic traits
follows a Brownian motion model in Euclidean space, in which the trait of each descendant
is derived from the trait of the ancestral species through a random walk in Euclidean space.
The time of this walk is proportional to the divergence time, which is the branch length of
the evolutionary tree.

% As demonstrated by \textcite{Wang2017NormFace}, in the last fully connected layer of CNNs,
% the semantic information is primarily encoded in the angular direction, which is equivalent
% to $L_2$ normalised vectors. Therefore, we assume that the feature vectors randomly walk on
% a unit hypersphere.
We implemented the algorithm for Spherical Ancestral State Reconstruction by modeling
phenotypic evolution as Riemannian Brownian Motion on the hypersphere of
the deep feature space in Python.

\begin{figure}[htbp]
    \centering
    \begin{minipage}{0.45\textwidth}
        \centering
        \begin{tikzpicture}[>=Stealth, thick, scale=0.9]

            % 1. Define Coordinates
            \coordinate (O) at (0,0);
            \coordinate (A_dir) at (100:6); % Direction of Va
            \coordinate (B_dir) at (10:6);  % Direction of Vb

            % Define the main radius
            \def\R{5}
            \coordinate (A) at (100:\R);
            \coordinate (B) at (10:\R);

            % Define P at an interpolated angle
            % Let's assume t is roughly 0.4 for visual similarity
            \coordinate (P) at (55:\R);

            % 2. Calculate Intersections for Parallelogram
            % We need A' on OA such that PA' // OB
            % We need B' on OB such that PB' // OA

            % Path for OA and OB
            \path[name path=rayOA] (O) -- (A_dir);
            \path[name path=rayOB] (O) -- (B_dir);

            % Path from P parallel to OB (to find A')
            \path[name path=lineFromP_parOB] (P) -- ++($(O)-(B_dir)$) -- ++($(B_dir)-(O)$);

            % Path from P parallel to OA (to find B')
            \path[name path=lineFromP_parOA] (P) -- ++($(O)-(A_dir)$) -- ++($(A_dir)-(O)$);

            % Find intersections
            \path [name intersections={of=rayOA and lineFromP_parOB, by=A_prime}];
            \path [name intersections={of=rayOB and lineFromP_parOA, by=B_prime}];

            % 3. Draw Geometry

            % Draw the main arc
            \draw (A) arc (100:10:\R);

            % Draw Vectors OA, OB, OP
            \draw[->] (O) -- (A) node[below left] {$\mathbf{v}_a$};
            \draw[->] (O) -- (B) node[below left] {$\mathbf{v}_b$};
            \draw[->] (O) -- (P) node[below left, xshift=-6pt, yshift=-2pt] {$\mathbf{v}_p$};

            % Draw Components (A' and B')
            % Vectors Va' and Vb' lie on the axes
            \draw[->, ultra thick, blue!80!black] (O) -- (A_prime) node[midway, left] {$\mathbf{v}_a'$};
            \draw[->, ultra thick, blue!80!black] (O) -- (B_prime) node[midway, below] {$\mathbf{v}_b'$};

            % Draw Dashed Parallel Lines
            \draw[dashed] (P) -- (A_prime);
            \draw[dashed] (P) -- (B_prime);

            % 4. Labels and Angles

            % Points labels
            \node[left] at (A_prime) {$A'$};
            \node[above right] at (B_prime) {$B'$};
            \node[above] at (A) {$A$};
            \node[right] at (B) {$B$};
            \node[above right] at (P) {$P$};

            % Angle arcs

            % Angle (1-t)theta between OA and OP
            \pic [draw, angle radius=1.0cm, pic text={$(1-t)\theta$}, angle eccentricity=1.6, font=\footnotesize] {angle = P--O--A};

            % Angle t*theta between OP and OB
            \pic [draw, angle radius=0.8cm, pic text={$t\theta$}, angle eccentricity=1.6, font=\footnotesize] {angle = B--O--P};

            % Internal angles (pi - theta) at A' and B'
            % Using small arcs inside the parallelogram
            \pic [draw, angle radius=0.6cm, pic text={$\pi-\theta$}, angle eccentricity=1.8, font=\footnotesize] {angle = O--A_prime--P};
            \pic [draw, angle radius=0.6cm, pic text={$\pi-\theta$}, angle eccentricity=1.8, font=\footnotesize] {angle = P--B_prime--O};

            % Angle t*theta at P (alternate interior) - optional based on diagram complexity
            % \pic [draw, angle radius=1.0cm, "$t\theta$"] {angle = A_prime--P--O};

        \end{tikzpicture}
    \end{minipage}
    \hfill
    \begin{minipage}{0.45\textwidth}
        \centering
        \begin{tikzpicture}[scale=2,every node/.style={font=\small},line width=0.8pt]
    \coordinate (Root) at (0,0);
    \coordinate (Internal) at (0, -0.3);
    \coordinate (A) at (-0.5, -2);
    \coordinate (B) at (0.5, -2);

    \draw (Root) -- (Internal);
    \draw (Internal) -| (A) node[pos=0.7, left] {$l_a$};
    \draw (Internal) -| (B) node[pos=0.7, right] {$l_b$};

    \node[below] at (A) {A};
    \node[below] at (B) {B};
    \node[above right] at (Internal) {P};
        \end{tikzpicture}
    \end{minipage}
    \caption{Geometric diagram of the ancestor state reconstruction algorithm. }
    \label{fig:algorithm-asr}
\end{figure}

In a post-order traversal of the phylogenetic tree, for a selected pair
of sister nodes (sister groups), let $\mathbf{v}_a, \mathbf{v}_b$
be the state vectors of them on a unit hypersphere, and $l_a, l_b$ be their branch
lengths (for leaf nodes) or the equivalent branch lengths (for internal nodes)
on a timetree. Let $\mathbf{v}_p$ be the feature vector of their parent node
(most recent common ancestor), it must be on the two-dimensional Euclidean plane spanned
by $\mathbf{v}_a$ and $\mathbf{v}_b$. The ancestral state and the contrast is computed
via the following steps \cref{fig:algorithm-asr}:

\begin{enumerate}
    \item We firstly calculate their geodesic distance, which is equivalent to their included angle
            for $L_2$ normalised vectors:
    \begin{equation}
        \theta = \arccos{(\mathbf{v}_a\cdot\mathbf{v}_b)}
    \end{equation}

    \item The parent state $\mathbf{v}_p$ is inferred by interpolating along the geodesic arc:
    \begin{equation}
        \mathbf{v}_p = \frac{\sin((1-t)\theta)}{\sin\theta}\mathbf{v}_a + \frac{\sin(t\theta)}{\sin\theta}\mathbf{v}_b
    \end{equation}
    where $t = l_b / (l_a + l_b)$.

    \item To adjust for the curvature of the state space, the contrast variance is scaled by the ratio of the squared arc length to the squared chord length:
    \begin{equation}
        \text{Contrast}^2 = \underbrace{\frac{(\mathbf{v}_a - \mathbf{v}_b)^{\circ 2}}{l_a + l_b}}_{\text{Euclidean Variance}} \times \underbrace{\frac{\theta^2}{||\mathbf{v}_a - \mathbf{v}_b||^2}}_{\text{Correction Factor}}
    \end{equation}
    where $\circ$ means Hadamard power (element-wise power). In standard Phylogenetic Comparative Methods (PCMs), the phenotypic distance between two states is measured as the linear chord length $||v_a - v_b||$. However, on a Riemannian manifold, the true evolutionary path follows the geodesic curve, with distance measured by the arc length $\theta$. For closely related species, the chord length approximates the arc length. However, for highly divergent taxa across deep evolutionary time, the linear chord length severely underestimates the true evolutionary distance. 

    \item When the traversal is completed, the estimated evolutionary rate of the whole tree is:
    \begin{equation}
        \hat{\sigma}^2 = \frac{\sum\text{Contrast}^2}{N-1}
    \end{equation}
    where $N$ is the numbers of leaf nodes (species).
\end{enumerate}

Once the whole ASR process is finished, the spherical variance for every time slice (1 ma) is calculated
on the whole timetree, which represents the phenotypic disparity of all living branches at the specific time.
Then, we conduct 100 times of null simulations based on Brownian motion model in preorder traversals.
For every internal node, the feature vectors of its children are calculated
according to the following algorithm (\cref{fig:algorithm-bm}):

\begin{figure}[htbp]
    \centering
    \begin{tikzpicture}[>=Stealth, thick, scale=0.7]

        % --- CONFIGURATION ---
        \def\R{6}          % Radius of the sphere (length of vp)
        \def\angTheta{45}  % Angle theta
        \def\angAlpha{65}  % Angle of vx relative to vertical
        \def\lenVx{4.5}    % Length of vx vector

        % Coordinates
        \coordinate (Origin) at (0,0);           % The "North Pole" (Tangent point)
        \coordinate (Center) at (0, -\R);        % The Center of the sphere

        % --- AXES ---
        % Vertical dashed axis (upwards)
        \draw[dashed, thick] (0, 0) -- (0, 2.5);

        % Horizontal dashed axis (Tangent Space)
        \draw[dashed, thick] (-2, 0) -- (7, 0) node[right] {tangent space};


        % --- TOP PART: PROJECTION ---

        % Vector vx
        \coordinate (Vx) at (\angAlpha:\lenVx); % Defined by angle alpha
        % Note: In the image, vx is (x,y). Let's approximate the slope visually.
        \coordinate (Vx_End) at (5, 2);
        \draw[->, ultra thick] (Origin) -- (Vx_End) node[midway, above left] {$\mathrm{v_x}$};

        % Vector vt (Projection on horizontal)
        \coordinate (Vt_End) at (5, 0);
        \draw[->, ultra thick] (Origin) -- (Vt_End) node[midway, above] {$\mathrm{v_t}$};

        % Projection line 'a'
        \draw[dashed, thick] (Vx_End) -- (Vt_End) node[midway, right] {$\mathrm{a}$};

        % Angle Alpha
        \coordinate (Y_Axis_Point) at (0, 2);
        \pic [draw, pic text=$\alpha$, angle radius=0.8cm, angle eccentricity=1.3] {angle = Vx_End--Origin--Y_Axis_Point};


        % --- BOTTOM PART: GEODESIC UPDATE ---

        % Calculate the "Child" point on the sphere surface
        % x = R * sin(theta)
        % y = -R + R * cos(theta) (relative to origin)
        % But we draw from Center (0, -R).
        % Target point S relative to Center is (R sin theta, R cos theta).

        \coordinate (S) at ({\R*sin(\angTheta)}, {-\R + \R*cos(\angTheta)});
        \coordinate (M) at (0, {-\R + \R*cos(\angTheta)}); % The corner of the triangle

        % 1. Vertical Component (vp * cos theta)
        % Draws from Center to M
        \draw[->, ultra thick] (Center) -- (M) node[midway, left=5pt] {$\mathrm{v_p} \cdot \cos\theta$};

        % 2. Horizontal Component (vt/|vt| * sin theta)
        % Draws from M to S
        \draw[->, ultra thick] (M) -- (S) node[midway, above] {$\frac{\mathrm{v_t}}{\|\mathrm{v_t}\|} \sin\theta$};

        % 3. Result Vector (vc)
        % Draws from Center to S
        \draw[->, ultra thick] (Center) -- (S) node[midway, below right] {$\mathrm{v_c}$};

        % 4. Parent Vector (vp) label logic
        % The line goes from Center to Origin.
        % We need an arrow head at the Origin.
        \draw[->, ultra thick] (Center) -- (Origin); % completing the vertical line
        \node[above=40pt, left=5pt] at (0, -2.5) {$\mathrm{v_p}$}; % Label roughly in the middle of the top section

        % 5. Angle Theta
        \pic [draw, pic text=$\theta$, angle radius=1.0cm, angle eccentricity=1.3] {angle = S--Center--M};

        % 6. The Geodesic Arc
        % Arcs from Origin (top) down to S.
        % To draw an arc in TikZ, we specify start angle, end angle, and radius.
        % Start: 90 degrees (relative to Center). End: 90 - theta.
        \draw[thick] (0,0) arc (90:{90-\angTheta}:\R);
    \end{tikzpicture}
    \caption{Geometric diagram of the brownian motion simulation algorithm. }
    \label{fig:algorithm-bm}
\end{figure}

\begin{enumerate}
    \item Create a Gaussian noise $\mathbf{v}_{x}$ in the Euclidean space following this multivariate normal distribution:
    \begin{equation}
        \mathbf{v}_{x} \sim \mathcal{N}(\mathbf{0}, \hat{\sigma}^2 t \mathbf{I}_D)
    \end{equation}
    where $t$ is the evolution time (branch length) of the child node.
    \item Project $\mathbf{v}_{x}$ onto the tangent space of the hypersphere:
    \begin{equation}
        a = \mathbf{v}_{p} \cdot \mathbf{v}_{x}
    \end{equation}
    \begin{equation}
        \mathbf{v}_{t} = \mathbf{v}_{x} - a \mathbf{v}_{p}
    \end{equation}
    \item Move on the hypersphere following a great circle, the direction and geodesic distance is same as the direction and the norm of $\mathbf{v}_{t}$:
    \begin{equation}
        \theta = ||\mathbf{v}_{t}||
    \end{equation}
    \begin{equation}
        \mathbf{v}_{c} = \cos{\theta}~\mathbf{v}_{p} + \frac{\sin{\theta}}{||\mathbf{v}_{t}||} \mathbf{v}_{t}
    \end{equation}
    and $\mathbf{v}_{c}$ is the simulational feature vector of the child species.
\end{enumerate}

When the whole traversal is finished, the ASR is conducted based on the simulational
feature vectors of all modern species (leaf nodes) following the aforementioned algorithm,
as well as the spherical variance for every time slice. Finally, the empirical disparity
through time result and the mean disparity through time of all 100 times of null simulations
are visualised and compared. The whole disparity-through-time analysis is conducted on the trees
by \textcite{Stiller2024Complexity}.